\section{Baseline Quasi-One-Dimensional Inviscid Nozzle Flow Model}
\label{sec:3}
% =========================================================

\subsection{Thermodynamic closure and governing relations}

\subsubsection{Governing assumptions}

The baseline analysis is based on the classical quasi-one-dimensional nozzle model. The following assumptions are used in the ideal formulation:
\begin{itemize}
    \item the flow is steady;
    \item the flow is adiabatic;
    \item the flow is inviscid in the core region;
    \item the nozzle flow is shock-free;
    \item the gas expansion is treated as isentropic;
    \item the flow is quasi-one-dimensional, meaning that properties vary mainly along the nozzle axis;
    \item the gas is approximated as a perfect gas with an effective specific heat ratio $\gamma$;
    \item sonic conditions occur at the throat for a choked nozzle;
    \item the throat area is equal to the critical area, $A_t=A^*$;
    \item the subsonic branch of the area--Mach relation is used upstream of the throat;
    \item the supersonic branch is used downstream of the throat.
\end{itemize}

These assumptions are appropriate for preliminary design. They are not sufficient for final qualification of a rocket nozzle because they neglect real-gas effects, finite-rate chemistry, wall heat transfer, boundary-layer separation, shocks, nonuniform inlet conditions, two-phase particles or droplets, and multidimensional wave structures.

\subsubsection{Thermochemical closure using NASA CEA/RocketCEA}

The thermochemical properties used by the simulator are obtained from NASA CEA through RocketCEA. For each combination of chamber pressure, mixture ratio and expansion ratio, RocketCEA can provide useful quantities such as combustion temperature, effective ratio of specific heats, molar mass of combustion products, characteristic velocity, thrust coefficient, exit Mach number, exit temperature, exit pressure and ideal specific impulse.

The classical relations below are written for a locally constant effective $\gamma$. In the implemented axial reconstruction, RocketCEA chamber, throat and exit values are interpolated along the contour and the local interpolated value is used in the area--Mach and stagnation-to-static relations. This remains an engineering property closure rather than a finite-rate or station-by-station equilibrium chemistry solution.

For a hybrid engine using nitrous oxide and paraffin-based fuel, RocketCEA also allows the user to define a custom fuel composition and evaluate equilibrium combustion properties for the chosen oxidizer-to-fuel ratio.

\subsubsection{Choked mass-flow constraint and external throat pre-sizing}

For a choked nozzle, the mass flow rate through the throat is governed by the stagnation conditions upstream of the throat and by the throat area. Assuming isentropic perfect-gas flow, the choked mass flow rate is
\begin{equation}
\dot{m}=A_t P_0
\sqrt{\frac{\gamma}{R T_0}}
\left(\frac{2}{\gamma+1}\right)^{\frac{\gamma+1}{2(\gamma-1)}} .
\label{eq:choked_mdot}
\end{equation}

Solving for throat area gives
\begin{equation}
A_t =
\frac{\dot{m}}{P_0
\sqrt{\dfrac{\gamma}{R T_0}}
\left(\dfrac{2}{\gamma+1}\right)^{\frac{\gamma+1}{2(\gamma-1)}}} ,
\label{eq:throat_area}
\end{equation}
with
\begin{equation}
r_t=\sqrt{\frac{A_t}{\pi}} .
\label{eq:throat_radius}
\end{equation}

This relation links the propulsion operating point to the geometry. For a fixed chamber pressure and combustion temperature, increasing mass flow requires a larger throat area. For a fixed mass flow rate, increasing chamber pressure allows the same mass flow through a smaller throat. In the present software, however, this calculation belongs to the upstream engine pre-sizing step: $r_t$ is supplied by the user and is neither recalculated nor optimized by the nozzle genetic algorithm.

\begin{figure}[H]
\centering
\begin{tikzpicture}[
    >=Latex,
    font=\small,
    station/.style={
        draw,
        rectangle,
        minimum width=2.8cm,
        minimum height=0.9cm,
        align=left,
        fill=white
    },
    box/.style={
        draw,
        rectangle,
        align=left,
        fill=white,
        inner sep=5pt
    },
    flowarrow/.style={
        ->,
        line width=0.8pt
    },
    leader/.style={
        gray,
        line width=0.5pt
    }
]

% =========================================================
% CEA block
% =========================================================
\node[align=left, anchor=north west] (cea) at (-3.0,4.0) {
\textbf{With CEA:}\\
-- Molar Mass $(M)$\\
-- Chamber Temperature $(T_c)$\\
-- Chamber Pressure $(P_c)$\\
-- Exit Temperature $(T_e)$\\
-- Exit Pressure $(P_e)$\\[2pt]
$R_u = 8.314~\mathrm{J/(mol\,K)}$\\
$R = R_u/M$
};

% =========================================================
% Formula block
% =========================================================
\node[box, minimum width=4.6cm, minimum height=1.4cm, anchor=north west] (formula) at (-3.75,0.3) {
$\displaystyle
A_{t,\mathrm{nozzle}}=
\frac{\dot{m}_{\mathrm{prop}}}
{P_c
\sqrt{\dfrac{\gamma_c}{R T_c}
\left(\dfrac{2}{\gamma_c+1}\right)^{\frac{\gamma_c+1}{\gamma_c-1}}}}
$
};

% =========================================================
% Bell-nozzle contour: same piecewise formulation as the 2D plot
% =========================================================

% ----------------------------
% Physical/geometry parameters used only for the schematic contour
% ----------------------------
\pgfmathsetmacro{\rt}{0.020}
\pgfmathsetmacro{\eps}{5.0}
\pgfmathsetmacro{\re}{sqrt(\eps)*\rt}

\pgfmathsetmacro{\alphaDeg}{15}
\pgfmathsetmacro{\bell}{0.80}
\pgfmathsetmacro{\thetaSubDeg}{60}
\pgfmathsetmacro{\thetaInDeg}{30}
\pgfmathsetmacro{\Rch}{0.060}
\pgfmathsetmacro{\tx}{0.055}

% ----------------------------
% Subsonic line and arc
% ----------------------------
\pgfmathsetmacro{\xr}{-1.5*\rt*sin(\thetaSubDeg)}
\pgfmathsetmacro{\yr}{2.5*\rt - 1.5*\rt*cos(\thetaSubDeg)}
\pgfmathsetmacro{\xR}{\tx+\xr}

\pgfmathsetmacro{\mLine}{-tan(\thetaSubDeg)}
\pgfmathsetmacro{\bLine}{\yr - \mLine*\xr}
\pgfmathsetmacro{\xf}{(\Rch - \bLine)/\mLine}
\pgfmathsetmacro{\xF}{\tx+\xf}

% ----------------------------
% Supersonic arc and parabola
% ----------------------------
\pgfmathsetmacro{\rSup}{0.4*\rt}
\pgfmathsetmacro{\Px}{\rSup*sin(\thetaInDeg)}
\pgfmathsetmacro{\Py}{1.4*\rt - \rSup*cos(\thetaInDeg)}
\pgfmathsetmacro{\xP}{\tx+\Px}

\pgfmathsetmacro{\Lcone}{(\re-\rt)/tan(\alphaDeg)}
\pgfmathsetmacro{\L}{\Px + \bell*\Lcone}
\pgfmathsetmacro{\xE}{\tx+\L}

\pgfmathsetmacro{\mZero}{tan(\thetaInDeg)}
\pgfmathsetmacro{\den}{(\L-\Px)}
\pgfmathsetmacro{\aPar}{(\re - \Py - \mZero*(\den))/(\den*\den)}
\pgfmathsetmacro{\bPar}{\mZero - 2*\aPar*\Px}
\pgfmathsetmacro{\cPar}{\Py - \aPar*\Px*\Px - \bPar*\Px}

% ----------------------------
% Drawing transformation
% These values only control the size and position of the schematic
% ----------------------------
\pgfmathsetmacro{\Xstart}{3.55}
\pgfmathsetmacro{\Xscale}{45.0}
\pgfmathsetmacro{\Yaxis}{1.55}
\pgfmathsetmacro{\Rscale}{20}

% Helper idea:
% plotted x = Xstart + Xscale*(x - xF)
% plotted upper y = Yaxis + Rscale*r(x)
% plotted lower y = Yaxis - Rscale*r(x)

% ----------------------------
% Upper wall
% ----------------------------

% 1) initial affine converging line
\draw[line width=0.95pt]
plot[domain=\xF:\xR, samples=2, variable=\x]
(
{\Xstart+\Xscale*(\x-\xF)},
{\Yaxis+\Rscale*(\mLine*(\x-\tx)+\bLine)}
);

% 2) subsonic circular throat arc
\draw[line width=0.95pt]
plot[domain=\xR:\tx, samples=120, variable=\x]
(
{\Xstart+\Xscale*(\x-\xF)},
{\Yaxis+\Rscale*(2.5*\rt - sqrt((1.5*\rt)^2 - (\x-\tx)^2))}
);

% 3) supersonic circular throat arc
\draw[line width=0.95pt]
plot[domain=\tx:\xP, samples=120, variable=\x]
(
{\Xstart+\Xscale*(\x-\xF)},
{\Yaxis+\Rscale*(1.4*\rt - sqrt((0.4*\rt)^2 - (\x-\tx)^2))}
);

% 4) parabolic divergent contour
\draw[line width=0.95pt]
plot[domain=\xP:\xE, samples=220, variable=\x]
(
{\Xstart+\Xscale*(\x-\xF)},
{\Yaxis+\Rscale*(\aPar*(\x-\tx)^2+\bPar*(\x-\tx)+\cPar)}
);

% ----------------------------
% Lower wall: axisymmetric mirror
% ----------------------------

% 1) initial affine converging line
\draw[line width=0.95pt]
plot[domain=\xF:\xR, samples=2, variable=\x]
(
{\Xstart+\Xscale*(\x-\xF)},
{\Yaxis-\Rscale*(\mLine*(\x-\tx)+\bLine)}
);

% 2) subsonic circular throat arc
\draw[line width=0.95pt]
plot[domain=\xR:\tx, samples=120, variable=\x]
(
{\Xstart+\Xscale*(\x-\xF)},
{\Yaxis-\Rscale*(2.5*\rt - sqrt((1.5*\rt)^2 - (\x-\tx)^2))}
);

% 3) supersonic circular throat arc
\draw[line width=0.95pt]
plot[domain=\tx:\xP, samples=120, variable=\x]
(
{\Xstart+\Xscale*(\x-\xF)},
{\Yaxis-\Rscale*(1.4*\rt - sqrt((0.4*\rt)^2 - (\x-\tx)^2))}
);

% 4) parabolic divergent contour
\draw[line width=0.95pt]
plot[domain=\xP:\xE, samples=220, variable=\x]
(
{\Xstart+\Xscale*(\x-\xF)},
{\Yaxis-\Rscale*(\aPar*(\x-\tx)^2+\bPar*(\x-\tx)+\cPar)}
);

% ----------------------------
% Throat marker
% ----------------------------
\draw[dashed, gray]
({\Xstart+\Xscale*(\tx-\xF)},{\Yaxis-\Rscale*\rt-0.04})
--
({\Xstart+\Xscale*(\tx-\xF)},{\Yaxis+\Rscale*\rt+0.04});

% =========================================================
% Station boxes
% =========================================================
\node (p1) at (2.7,1.55) {
\begin{tabular}{c}
$P_c$\\
$T_c$
\end{tabular}
};

\node (p2) at (9.65,1.5) {
\begin{tabular}{c}
$P_e$\\
$T_e$
\end{tabular}
};

% =========================================================
% Mass flow / gamma box
% =========================================================
\node[box,font=\footnotesize,
      minimum width=2.9cm,
      minimum height=0.7cm,
      inner sep=2pt] (massbox) at (6.70,0.23) {
Propellat Mass Flow $(\dot{m}_{\text{prop}})$\\
Specific Heat Ratio $(\gamma_c)$
};

% =========================================================
% Flow arrows
% =========================================================
\draw[flowarrow] (7.55,1.55) -- (9.25,1.55);
\draw[flowarrow] (7.65,1.82) -- (9.05,1.82);
\draw[flowarrow] (7.65,1.28) -- (9.05,1.28);

\draw[flowarrow] (3.15,1.74) -- (4.5,1.74);
\draw[flowarrow] (3.15,1.36) -- (4.5,1.36);

% small flow streak lines
\draw[leader] (7.25,1.70) -- (7.95,1.70);
\draw[leader] (7.25,1.40) -- (7.95,1.40);

\end{tikzpicture}
\caption{Nozzle flow stations used for external throat pre-sizing and ideal one-dimensional analysis.}
\label{fig:flow_stations_tikz}
\end{figure}

\subsubsection{Thrust balance and ideal expansion criterion}

The thrust produced by a rocket nozzle can be approximated by
\begin{equation}
F=\dot{m}V_e+(p_e-p_a)A_e,
\label{eq:thrust}
\end{equation}
where $V_e$ is the exit velocity, $p_e$ is the nozzle exit pressure, $p_a$ is the ambient pressure and $A_e$ is the exit area.

The first term is the momentum thrust. The second term is the pressure thrust. For a fixed chamber condition and mass flow rate, the ideal expansion condition at a given altitude is approximately
\begin{equation}
p_e=p_a .
\label{eq:ideal_expansion}
\end{equation}

When $p_e>p_a$, the nozzle is underexpanded and the gases still have remaining expansion potential at the exit. When $p_e<p_a$, the nozzle is overexpanded, which can lead to compression waves, shocks, separation risk and loss of performance depending on the operating regime.

Because ambient pressure varies during flight, a fixed-geometry nozzle cannot be perfectly expanded at all altitudes. For low-altitude or short-duration flights, the simulator may use a representative ambient pressure, such as sea-level pressure or a trajectory-averaged ambient pressure. For a more complete flight analysis, the expansion ratio should be evaluated over the expected altitude range and the selected value should minimize a trajectory-weighted performance penalty.

\subsubsection{Expansion ratio and exit-plane radius}

The expansion ratio is defined as
\begin{equation}
\varepsilon=\frac{A_e}{A_t} .
\label{eq:expansion_ratio}
\end{equation}

For circular sections,
\begin{equation}
\varepsilon=\frac{\pi r_e^2}{\pi r_t^2}=\left(\frac{r_e}{r_t}\right)^2,
\end{equation}
therefore
\begin{equation}
r_e=\sqrt{\varepsilon}\,r_t .
\label{eq:exit_radius}
\end{equation}

The expansion ratio is one of the most important nozzle design variables. Increasing $\varepsilon$ generally increases exit Mach number and ideal exhaust velocity, but also increases nozzle length, exit area, structural mass and overexpansion risk at low altitude.

\subsubsection{Propellant mass-flow closure from mixture ratio}

Using the conventional definition
\begin{equation}
O/F=\frac{\dot{m}_{ox}}{\dot{m}_{fuel}},
\end{equation}
the fuel mass flow rate is
\begin{equation}
\dot{m}_{fuel}=\frac{\dot{m}_{ox}}{O/F},
\end{equation}
and the total propellant mass flow rate is
\begin{equation}
\dot{m}_{prop}=\dot{m}_{ox}+\dot{m}_{fuel}
=\dot{m}_{ox}\left(1+\frac{1}{O/F}\right).
\label{eq:mdot_total}
\end{equation}

This convention must be used consistently. Reversing the definition of mixture ratio changes the mass-flow calculation and therefore the throat sizing.

% 
% =========================================================

\subsection{Axial flow-field reconstruction}

\subsubsection{Stagnation-to-static isentropic relations}

For one-dimensional, steady, adiabatic and inviscid flow of a perfect gas, the stagnation-to-static relations are
\begin{align}
\frac{T}{T_0} &= \left(1+\frac{\gamma-1}{2}M^2\right)^{-1},
\label{eq:T_T0}\\
\frac{p}{p_0} &= \left(1+\frac{\gamma-1}{2}M^2\right)^{-\gamma/(\gamma-1)},
\label{eq:p_p0}\\
\frac{\rho}{\rho_0} &= \left(1+\frac{\gamma-1}{2}M^2\right)^{-1/(\gamma-1)}.
\label{eq:rho_rho0}
\end{align}

Once the Mach number distribution $M(x)$ is known, these relations provide the axial distributions of temperature, pressure and density.

\subsubsection{Area--Mach compatibility relation}

The Mach number is implicitly related to the local cross-sectional area through the area--Mach relation:
\begin{equation}
\frac{A}{A^*}=\frac{1}{M}
\left[
\frac{2}{\gamma+1}
\left(1+\frac{\gamma-1}{2}M^2\right)
\right]^{\frac{\gamma+1}{2(\gamma-1)}} .
\label{eq:area_mach}
\end{equation}

For a choked nozzle, $A^*=A_t$. Therefore, at each axial station,
\begin{equation}
\frac{A(x)}{A_t}=\frac{\pi r(x)^2}{\pi r_t^2}=\left(\frac{r(x)}{r_t}\right)^2 .
\label{eq:area_ratio_local}
\end{equation}

Equation~\eqref{eq:area_mach} has two possible Mach solutions for $A/A^*>1$: a subsonic solution and a supersonic solution. The correct branch is selected from the location relative to the throat. Upstream of the throat, the subsonic solution is used. Downstream of the throat, the supersonic solution is used. At the throat, $M=1$.

\subsubsection{Branch-resolved numerical solution of the Mach number}

Because Eq.~\eqref{eq:area_mach} cannot generally be inverted explicitly for $M$, the simulator determines $M(x)$ by solving
\begin{equation}
f(M)=\frac{A}{A^*} - \frac{1}{M}
\left[
\frac{2}{\gamma+1}
\left(1+\frac{\gamma-1}{2}M^2\right)
\right]^{\frac{\gamma+1}{2(\gamma-1)}}=0 .
\label{eq:area_mach_residual}
\end{equation}

The implementation uses Brent's bracketed root method, combining the robustness of a bracketing method with faster interpolation steps when the residual permits them. The numerical intervals are
\begin{equation}
M\in(0,1) \quad \text{for the subsonic branch},
\end{equation}
\begin{equation}
M\in(1,M_{max}) \quad \text{for the supersonic branch}.
\end{equation}

Once $M(x)$ is obtained, the local speed of sound and axial velocity are
\begin{equation}
a(x)=\sqrt{\gamma R T(x)},
\end{equation}
\begin{equation}
u(x)=M(x)a(x).
\label{eq:velocity}
\end{equation}


\subsubsection{Axial reconstruction of Mach, static-pressure and static-temperature curves}

The axial property curves couple the piecewise radius function to the
quasi-one-dimensional isentropic model. The process starts from the discretized
nozzle contour,
\begin{equation}
\left\{x_i,r_i\right\}_{i=0}^{N},
\end{equation}
where each station has the local cross-sectional area
\begin{equation}
A_i=\pi r_i^2.
\label{eq:local_area_discrete}
\end{equation}

Because the throat is assumed choked, the critical area is
\begin{equation}
A^*=A_t=\pi r_t^2.
\end{equation}
Therefore, the local area ratio used in the area--Mach equation is
\begin{equation}
\Lambda_i=\frac{A_i}{A_t}=\left(\frac{r_i}{r_t}\right)^2.
\label{eq:local_area_ratio_discrete}
\end{equation}

For each station, the Mach number is obtained by solving the residual form of the area--Mach equation,
\begin{equation}
F(M_i,\Lambda_i)=
\frac{1}{M_i}
\left[
\frac{2}{\gamma+1}
\left(1+\frac{\gamma-1}{2}M_i^2\right)
\right]^{\frac{\gamma+1}{2(\gamma-1)}}
-\Lambda_i=0.
\label{eq:mach_curve_residual}
\end{equation}

The solution branch is selected from the geometric station location:
\begin{equation}
0<M_i<1 \quad \text{for} \quad x_i<x_t,
\end{equation}
\begin{equation}
M_i=1 \quad \text{for} \quad x_i=x_t,
\end{equation}
\begin{equation}
M_i>1 \quad \text{for} \quad x_i>x_t.
\end{equation}
This branch selection is essential because the area--Mach relation is double-valued for \(A/A^*>1\). The same area ratio may correspond either to a subsonic state upstream of the throat or to a supersonic state downstream of it.

After the Mach number has been obtained at each station, the static-temperature curve is generated directly from the stagnation-to-static relation,
\begin{equation}
T_i=T(x_i)=\frac{T_0}{1+\dfrac{\gamma-1}{2}M_i^2}.
\label{eq:temperature_curve}
\end{equation}
Similarly, the static-pressure curve is obtained from
\begin{equation}
p_i=p(x_i)=p_0
\left(1+\frac{\gamma-1}{2}M_i^2\right)^{-\gamma/(\gamma-1)}.
\label{eq:pressure_curve}
\end{equation}
The density curve may also be reconstructed as
\begin{equation}
\rho_i=\rho(x_i)=\rho_0
\left(1+\frac{\gamma-1}{2}M_i^2\right)^{-1/(\gamma-1)}.
\label{eq:density_curve}
\end{equation}

The simulator therefore stores the axial flow-field dataset as
\begin{equation}
\mathcal{D}_{flow}=\left\{x_i,\,A_i,\,\Lambda_i,\,M_i,\,T_i,\,p_i,\,\rho_i,\,u_i\right\}_{i=0}^{N},
\label{eq:flow_dataset}
\end{equation}
where the velocity is obtained from
\begin{equation}
u_i=M_i\sqrt{\gamma R T_i}.
\label{eq:velocity_curve}
\end{equation}

The Mach, pressure and temperature curves are then plotted as functions of the translated axial coordinate \(\bar{x}\). This means that the plotted distributions are not read directly from CEA along the nozzle. CEA/RocketCEA provides the thermochemical closure and reference stagnation properties, while the spatial evolution along the generated contour is reconstructed from \(r(x)\), the area--Mach relation and the isentropic equations.

In the converging subsonic region, reducing area accelerates the flow and increases Mach number toward unity. At the throat, the nozzle becomes choked and $M=1$. In the diverging supersonic region, increasing area further accelerates the flow and produces a supersonic Mach number increase.

The temperature and pressure decrease throughout the expansion according to Eqs.~\eqref{eq:T_T0} and~\eqref{eq:p_p0}. Therefore, the simulator can generate axial distributions of Mach number, pressure and temperature from the geometric contour alone, once the chamber stagnation state and thermodynamic model have been selected.
